% Part: first-order-logic
% Chapter: syntax
% Section: intro-syntax

\documentclass[../../../include/open-logic-section]{subfiles}

\begin{document}

\olfileid{fol}{syn}{itx}

\olsection{પરિચય}

પ્રથમ-ક્રમ તર્કશાસ્ત્રના સિદ્ધાંત અને અધિસિદ્ધાંતનો વિકાસ કરવા માટે
આપણે પહેલાં તેની અભિવ્યક્તિઓની વાક્યરચના અને અર્થવિચાર વ્યાખ્યાયિત
કરવાં પડે. પ્રથમ-ક્રમ તર્કશાસ્ત્રની અભિવ્યક્તિઓ પદો અને
!!{formula}s છે. પદો !!{variable}s, !!{constant}s અને !!{function}sમાંથી
બને છે. પછી !!^{formula}s, !!{predicate}s સાથે પદો જોડીને રચાય છે
(તેનાથી સૌથી નાનાં, ``આણ્વિક'' !!{formula}s બને છે); અને આણ્વિક
!!{formula}sમાંથી તાર્કિક સંયોજકો અને પરિમાણકો વડે વધુ જટિલ સૂત્રો
રચી શકીએ છીએ. રચના-નિયમો આપવાની ઘણી જુદી રીતો છે; અહીં આપણે માત્ર
એક શક્ય રીત આપીએ છીએ. બીજાં તંત્રો જુદા સંકેતો પસંદ કરશે, જુદા
સંયોજક-ગણોને મૂળભૂત ગણશે અને કૌંસનો જુદી રીતે ઉપયોગ કરશે (અથવા
કહેવાતી પોલિશ સંકેતપદ્ધતિની જેમ કૌંસ બિલકુલ નહિ વાપરે). જોકે બધા
અભિગમોમાં સમાન વાત એ છે કે રચના-નિયમો પદો અને !!{formula}sના ગણને
\emph{અનુમાનાત્મક રીતે} વ્યાખ્યાયિત કરે છે. આ યોગ્ય રીતે કરવામાં આવે
તો દરેક અભિવ્યક્તિ રચના-નિયમો મુજબ મૂળભૂત રીતે માત્ર એક જ રીતે બની
શકે છે. અભિવ્યક્તિઓને \emph{એકમાત્ર વાચનીય} બનાવતી આ
અનુમાનાત્મક વ્યાખ્યાને લીધે આપણે એ જ પદ્ધતિ---અનુમાનાત્મક
વ્યાખ્યા---વાપરીને આ અભિવ્યક્તિઓને અર્થ આપી શકીએ છીએ.

\end{document}
